\chapter{Type--Class Unification, Descriptors, and C3 MRO (Python 2.2--2.3)}

Python 2.2 made the single most important change to the object model after its inception: it
\textbf{unified types and classes}. Before it, user-defined classes and built-in types were
different kinds of thing at the C level, and you could not subclass \texttt{list}. After it, a
class \emph{is} a type, every class descends from \texttt{object}, and three mechanisms---the
\textbf{descriptor protocol}, the \textbf{C3 linearized MRO}, and the
\textbf{\_\_new\_\_/\_\_init\_\_ split}---became the machinery underneath properties, methods,
\texttt{super()}, and \texttt{@classmethod}. This chapter is the origin story and the working
reference for all three; it is where the dunder-to-slot wiring of Chapter 2 becomes a usable
protocol you can implement yourself.

\textbf{One home.} Descriptors and the MRO are taught in full here. Metaclasses and the
construction of \texttt{type} itself get their dedicated treatment in Vol VIII; this chapter
establishes the model they build on.

\section*{Section Index}
\begin{itemize}
  \item \textbf{4.1} New-style classes: unifying types and classes (PEP 252/253)
  \item \textbf{4.2} The descriptor protocol
  \item \textbf{4.3} The attribute-lookup algorithm and descriptor precedence
  \item \textbf{4.4} Methods, \texttt{classmethod}, \texttt{staticmethod}, and \texttt{property} as descriptors
  \item \textbf{4.5} Method Resolution Order and C3 linearization
  \item \textbf{4.6} Instance lifecycle: \texttt{\_\_new\_\_} vs \texttt{\_\_init\_\_}
  \item \textbf{4.7} Performance, the senior-engineer contrast, and anti-patterns
  \item \textbf{4.8} Summary and cross-references
\end{itemize}

\section{New-style classes: unifying types and classes (PEP 252/253)}

\textbf{Why this exists.} Before Python 2.2 there were two parallel universes. \emph{Classic
classes} (\texttt{class C:}) were all instances of one C type, \texttt{PyClass\_Type}; their
instances were all \texttt{PyInstance\_Type}. \emph{Built-in types} were
\texttt{PyTypeObject}s. The two could not mix: \texttt{type(c)} of a classic instance returned
\texttt{<type 'instance'>}, and \textbf{you could not subclass a built-in type}. \textbf{PEP
252} and \textbf{PEP 253} introduced \textbf{new-style classes}: a class deriving from
\texttt{object} is itself a \texttt{PyTypeObject}---an instance of \texttt{type}.

\begin{lstlisting}[language=Python, caption={In Python 3 every class is new-style; type(instance) IS the class.}]
class NewStyle:        # implicitly subclasses object in Python 3
    pass

n = NewStyle()
print("type(n):", type(n).__name__)
print("issubclass(NewStyle, object):", issubclass(NewStyle, object))
print("type(NewStyle) is type:", type(NewStyle) is type)
\end{lstlisting}

Verified output (CPython 3.13.5):

\begin{lstlisting}[caption={Output}]
type(n): NewStyle
issubclass(NewStyle, object): True
type(NewStyle) is type: True
\end{lstlisting}

The key C-level consequence (Chapter 2): every class is a \texttt{PyTypeObject} whose
\textbf{slots} hold the function pointers the interpreter calls. A Python \texttt{\_\_init\_\_}
installs a \textbf{slot wrapper} (\texttt{slot\_tp\_init}) into \texttt{tp\_init} that calls
your function; a built-in's C slot is exposed to Python as a \textbf{wrapper descriptor}. Python
methods and C slots are bridged in both directions.

\section{The descriptor protocol}

\textbf{Why this exists.} A \textbf{descriptor} customizes what happens when it is accessed as a
\emph{class attribute} of another object. It is the single mechanism behind methods,
\texttt{property}, \texttt{classmethod}, \texttt{staticmethod}, \texttt{\_\_slots\_\_} members,
and \texttt{cached\_property}. The protocol is three optional methods: \texttt{\_\_get\_\_(self,
instance, owner=None)}, \texttt{\_\_set\_\_(self, instance, value)}, and
\texttt{\_\_delete\_\_(self, instance)}.

The presence of \texttt{\_\_set\_\_}/\texttt{\_\_delete\_\_} splits descriptors into two kinds:
a \textbf{data descriptor} defines \texttt{\_\_set\_\_} and/or \texttt{\_\_delete\_\_}
(\texttt{property} is canonical); a \textbf{non-data descriptor} defines only
\texttt{\_\_get\_\_} (a plain function is canonical).

\begin{lstlisting}[language=Python, caption={A minimal data descriptor that validates writes.}]
class Positive:
    def __set_name__(self, owner, name):      # 3.6+: learn the attribute name
        self.private = "_" + name
    def __get__(self, obj, owner=None):
        if obj is None:
            return self                        # accessed on the class
        return getattr(obj, self.private)
    def __set__(self, obj, value):
        if value <= 0:
            raise ValueError("must be positive")
        setattr(obj, self.private, value)

class Account:
    balance = Positive()
    def __init__(self, b):
        self.balance = b

a = Account(100)
print("balance:", a.balance)
try:
    a.balance = -5
except ValueError as e:
    print("rejected:", e)
\end{lstlisting}

Verified output (CPython 3.13.5):

\begin{lstlisting}[caption={Output}]
balance: 100
rejected: must be positive
\end{lstlisting}

\texttt{\_\_set\_name\_\_} (PEP 487, Python 3.6) tells the descriptor its own attribute name at
class-creation time, removing 2.x-era boilerplate.

\section{The attribute-lookup algorithm and descriptor precedence}

\textbf{What the interpreter actually does.} Every \texttt{obj.name} read goes through the
type's \texttt{tp\_getattro} slot, which for ordinary objects is
\texttt{PyObject\_GenericGetAttr} (\texttt{Objects/object.c}):
\begin{enumerate}
  \item \textbf{Search the MRO} (\texttt{\_PyType\_Lookup}) for \texttt{name}; call any result
    \texttt{descr}.
  \item \textbf{Data-descriptor short-circuit.} If \texttt{descr} is a data descriptor, call
    \texttt{descr.\_\_get\_\_(obj, type(obj))} and return---\emph{the instance dict is never
    consulted}.
  \item \textbf{Instance dict.} Otherwise, if \texttt{name} is in \texttt{obj.\_\_dict\_\_},
    return it.
  \item \textbf{Non-data descriptor / class attribute.} Otherwise, if \texttt{descr} exists:
    if non-data, call \texttt{\_\_get\_\_}; if a plain class attribute, return it.
  \item \textbf{\_\_getattr\_\_ fallback.} Otherwise raise \texttt{AttributeError}, which
    triggers \texttt{\_\_getattr\_\_} if defined.
\end{enumerate}

\begin{lstlisting}[caption={Attribute-lookup flow}]
obj.name
  -> MRO has a DATA descriptor?  yes -> descr.__get__(obj, type)  (dict ignored)
  -> name in obj.__dict__?       yes -> return obj.__dict__[name]
  -> MRO has NON-DATA descriptor? yes -> descr.__get__(obj, type)
  -> MRO has plain class attr?   yes -> return it
  -> raise AttributeError -> __getattr__(self, name) if defined
\end{lstlisting}

\begin{lstlisting}[language=Python, caption={Data descriptors beat the instance dict; non-data descriptors do not.}]
class DataDesc:
    def __get__(self, obj, owner): return "data-descriptor"
    def __set__(self, obj, value): pass          # makes it a DATA descriptor
class NonDataDesc:
    def __get__(self, obj, owner): return "non-data-descriptor"

class C:
    d = DataDesc()
    n = NonDataDesc()

c = C()
c.__dict__["d"] = "instance-d"     # attempt to shadow the data descriptor
c.__dict__["n"] = "instance-n"     # attempt to shadow the non-data descriptor
print("c.d ->", c.d)
print("c.n ->", c.n)
\end{lstlisting}

Verified output (CPython 3.13.5):

\begin{lstlisting}[caption={Output}]
c.d -> data-descriptor
c.n -> instance-n
\end{lstlisting}

This is why you cannot accidentally clobber a \texttt{property} (data descriptor) by assigning
\texttt{self.x} in \texttt{\_\_init\_\_}, but you \emph{can} shadow a method (non-data
descriptor) with an instance attribute---and why ordinary instance attributes are cheap (they
resolve at step 3).

\section{Methods, \texttt{classmethod}, \texttt{staticmethod}, and \texttt{property} as descriptors}

Every flavor of ``method'' is a descriptor, differing only in \texttt{\_\_get\_\_}: a
\textbf{plain function} is a non-data descriptor whose \texttt{\_\_get\_\_} returns a
\textbf{bound method} pairing function and instance; \texttt{@classmethod} binds the
\textbf{class}; \texttt{@staticmethod} returns the function unchanged; \texttt{@property} is a
\textbf{data descriptor} wrapping \texttt{fget}/\texttt{fset}/\texttt{fdel}.

\begin{lstlisting}[language=Python, caption={A method is a function (non-data descriptor) that binds via \_\_get\_\_.}]
class Demo:
    def method(self):
        return "called"

d = Demo()
print("type(d.method).__name__:", type(d.method).__name__)       # method (bound)
print("type(Demo.method).__name__:", type(Demo.method).__name__) # function (raw)
print("manual __get__ == d.method:", Demo.method.__get__(d, Demo) == d.method)
\end{lstlisting}

Verified output (CPython 3.13.5):

\begin{lstlisting}[caption={Output}]
type(d.method).__name__: method
type(Demo.method).__name__: function
manual __get__ == d.method: True
\end{lstlisting}

The bound method holds \texttt{obj} in \texttt{\_\_self\_\_} and the function in
\texttt{\_\_func\_\_}; calling it prepends \texttt{\_\_self\_\_}. That is the entire mechanism
of \texttt{self}.

\section{Method Resolution Order and C3 linearization}

\textbf{Why this exists.} With multiple inheritance, ``which base defines this method?'' needs a
\emph{consistent} total order over ancestors. Classic depth-first, left-to-right reached a
shared ancestor before a sibling that overrode it. Python 2.3 adopted \textbf{C3
linearization}, guaranteeing \textbf{local precedence} (bases in declared order) and
\textbf{monotonicity} (if \texttt{X} precedes \texttt{Y} in one MRO, it does in every
subclass's MRO).

Let \texttt{L(C)} be the MRO of \texttt{C} with parents $B_1 \dots B_n$:
\[ L(C) = [C] + \mathrm{merge}\big(L(B_1), \dots, L(B_n), [B_1, \dots, B_n]\big) \]
\texttt{merge} repeatedly takes a \textbf{good head}---the head of some list absent from the
\emph{tail} of every other list---appends it, and removes it everywhere. If none exists, CPython
raises \texttt{TypeError}.

\textbf{Worked example (simple diamond).} For $O\to\{X,Y\}\to A$: $L(X)=[X,O,\text{object}]$,
$L(Y)=[Y,O,\text{object}]$, and merging yields $L(A)=[A,X,Y,O,\text{object}]$.

\begin{lstlisting}[language=Python, caption={C3 in practice --- confirm the worked result against the real \_\_mro\_\_.}]
class O: pass
class X(O): pass
class Y(O): pass
class A(X, Y): pass
print("A.__mro__:", [c.__name__ for c in A.__mro__])
\end{lstlisting}

Verified output (CPython 3.13.5):

\begin{lstlisting}[caption={Output}]
A.__mro__: ['A', 'X', 'Y', 'O', 'object']
\end{lstlisting}

\textbf{Worked example (Pedroni's monotonicity stress test).}

\begin{lstlisting}[language=Python, caption={The Pedroni example --- C3 produces a single consistent order.}]
class Ao: pass
class Bo: pass
class Co: pass
class Do: pass
class Eo: pass
class K1(Ao, Bo, Co): pass
class K2(Do, Bo, Eo): pass
class K3(Do, Ao): pass
class Z(K1, K2, K3): pass
print("Z.__mro__:", [c.__name__ for c in Z.__mro__])
\end{lstlisting}

Verified output (CPython 3.13.5):

\begin{lstlisting}[caption={Output}]
Z.__mro__: ['Z', 'K1', 'K2', 'K3', 'Do', 'Ao', 'Bo', 'Co', 'Eo', 'object']
\end{lstlisting}

This matches the hand computation $[Z, K1, K2, K3, D, A, B, C, E, \text{object}]$ exactly.

\begin{lstlisting}[language=Python, caption={An impossible hierarchy is rejected at class-creation time.}]
try:
    class Bad1: pass
    class Bad2(Bad1): pass
    class Bad3(Bad1, Bad2): pass     # demands Bad1 before AND after Bad2
except TypeError as e:
    print("TypeError:", e)
\end{lstlisting}

Verified output (CPython 3.13.5):

\begin{lstlisting}[caption={Output}]
TypeError: Cannot create a consistent method resolution order (MRO) for bases Bad1, Bad2
\end{lstlisting}

\textbf{\texttt{super()} follows the MRO}, dispatching to the next class in
\texttt{type(self).\_\_mro\_\_} after the current one---which is why cooperative multiple
inheritance works only when every class calls \texttt{super()}.

\section{Instance lifecycle: \texttt{\_\_new\_\_} vs \texttt{\_\_init\_\_}}

CPython splits construction into two slots: \textbf{\texttt{\_\_new\_\_}} ($\rightarrow$
\texttt{tp\_new}), a static method that \textbf{allocates} and returns the instance (ultimately
\texttt{PyType\_GenericNew}, setting up the \texttt{PyObject} header), running first; and
\textbf{\texttt{\_\_init\_\_}} ($\rightarrow$ \texttt{tp\_init}), an instance method that
\textbf{initializes} the object and returns \texttt{None}, running only if \texttt{\_\_new\_\_}
returned an instance of \texttt{cls}.

\begin{lstlisting}[language=Python, caption={\_\_new\_\_ allocates and runs first; \_\_init\_\_ initializes after.}]
class Lifecycle:
    def __new__(cls, *args, **kwargs):
        print("  __new__ (allocate)")
        return super().__new__(cls)
    def __init__(self, v):
        print("  __init__ (initialize) with", v)
        self.v = v

obj = Lifecycle(7)
print("obj.v:", obj.v)
\end{lstlisting}

Verified output (CPython 3.13.5):

\begin{lstlisting}[caption={Output}]
  __new__ (allocate)
  __init__ (initialize) with 7
obj.v: 7
\end{lstlisting}

This split is what makes immutable types possible: \texttt{int}, \texttt{str}, \texttt{tuple},
and frozen dataclasses do their work in \texttt{\_\_new\_\_}. Overriding \texttt{\_\_new\_\_} to
return an object of a \emph{different} type silently skips \texttt{\_\_init\_\_}.

\section{Performance, the senior-engineer contrast, and anti-patterns}

\textbf{The senior-engineer contrast.} C++ supports multiple inheritance via virtual
inheritance and per-class vtables, with no language-level linearization guarantee. Java forbids
multiple class inheritance (only interfaces). Python's C3 gives multiple inheritance a single,
predictable, monotonic order and a cooperative \texttt{super()}.

\textbf{Attribute-lookup cost.} The MRO walk of §4.3 would be expensive on every access, so
CPython caches \texttt{\_PyType\_Lookup} results keyed by a \textbf{type version tag}
invalidated on class mutation, and the 3.11+ specializing interpreter specializes
\texttt{LOAD\_ATTR} with inline caches (Vol VI). Practical rules: \texttt{\_\_slots\_\_}
(Vol VIII) removes the per-instance \texttt{\_\_dict\_\_} (less memory, faster access) and is
itself built from data descriptors; mutating classes at runtime invalidates the version cache
and de-optimizes dependent call sites.

\textbf{Anti-patterns.} Forgetting \texttt{super().\_\_init\_\_()} in a multiple-inheritance
hierarchy; implementing only \texttt{\_\_get\_\_} when you need write interception (use a data
descriptor); doing real work in \texttt{\_\_init\_\_} for an immutable type; heavy
\texttt{\_\_getattribute\_\_} (runs on every access; always delegate via
\texttt{super().\_\_getattribute\_\_}).

\section{Summary and cross-references}

\begin{itemize}
  \item Python 2.2 \textbf{unified types and classes}; Python 3 has only new-style classes.
  \item A \textbf{descriptor} customizes access via \texttt{\_\_get\_\_}/\texttt{\_\_set\_\_}/\texttt{\_\_delete\_\_};
    \textbf{data descriptors} outrank the instance dict, \textbf{non-data} do not---explaining
    methods, \texttt{property}, \texttt{classmethod}, \texttt{staticmethod}, \texttt{\_\_slots\_\_}.
  \item \texttt{PyObject\_GenericGetAttr} resolves: data descriptor $\rightarrow$ instance dict
    $\rightarrow$ non-data/class attr $\rightarrow$ \texttt{\_\_getattr\_\_}.
  \item \textbf{C3 linearization} gives multiple inheritance \textbf{local precedence} and
    \textbf{monotonicity}; \texttt{super()} walks it; inconsistent hierarchies raise
    \texttt{TypeError}.
  \item Construction is \textbf{\_\_new\_\_ (allocate) then \_\_init\_\_ (initialize)}.
  \item CPython makes lookup fast via the \textbf{type version cache} and 3.11+
    \texttt{LOAD\_ATTR} specialization; runtime class mutation defeats both.
\end{itemize}

\textbf{Cross-references.} \texttt{PyObject}/\texttt{PyTypeObject} and dunder-to-slot wiring
$\rightarrow$ Chapter 2. \texttt{property} and context managers $\rightarrow$ Chapter 5.
Metaclasses and the construction of \texttt{type} $\rightarrow$ Vol VIII. \texttt{\_\_slots\_\_}
and attribute-access optimization $\rightarrow$ Vol VIII. The specializing interpreter and
\texttt{LOAD\_ATTR} inline caches $\rightarrow$ Vol VI, Ch 17. \texttt{cached\_property}
$\rightarrow$ Vol IV.
